\chapter{Positive Matrices}\label{ch:perron}

\begin{center}\itshape
A certain man had put a pair of rabbits in a place surrounded on all sides by a wall. How many pairs of rabbits can be
bred from that pair in a year, if it is supposed that every month each pair begets a new pair, which from the second
month on becomes productive?
\upshape(Leonardo of Pisa, \emph{Liber abaci}, 1202)
\end{center}

In 1907 Oskar Perron, a young mathematician in Munich, was studying continued fractions and needed to know how the
powers of a matrix with positive entries behave. He proved that such a matrix has a positive eigenvalue strictly larger
in modulus than all the others, with an eigenvector whose entries are all positive. Georg Frobenius, then the grand old
man of Berlin algebra, read the paper, admired it, and spent the next five years extending it to matrices with
non-negative entries, a much subtler problem: there, several eigenvalues may share the largest modulus, and Frobenius
found that they are then arranged with perfect symmetry, like the vertices of a regular polygon.

The theorem lay quietly in algebra books for decades and then turned out to be everywhere. In 1906 Andrey Markov had
invented his chains to show that the law of large numbers does not need independence; their long-run behaviour is
Perron--Frobenius theory for stochastic matrices. In 1945 Patrick Leslie, studying the brown rat, wrote populations
structured by age as a non-negative matrix, whose Perron eigenvalue is the growth rate of the population and whose Perron
vector is its stable age distribution; Fibonacci's rabbits are the simplest case. Economists (Leontief's input--output
models) and demographers followed. And in 1998 two Stanford students, Sergey Brin and Larry Page, ranked the pages of the
web by the Perron vector of a huge stochastic matrix. A 2006 article in \emph{SIAM Review} called it ``the
\$25{,}000{,}000{,}000 eigenvector''.

This chapter proves Perron's theorem twice, once by an optimisation (Collatz--Wielandt) and once by a contraction that also
proves the convergence of the power method; extends it to irreducible matrices (Frobenius); studies periods and
primitivity, with Wielandt's sharp bound; and applies the results to Markov chains, PageRank and populations.

\begin{think}
Fibonacci's rabbits: write $(J_k,A_k)$ for the numbers of juvenile and adult pairs in month $k$. Then $J_{k+1}=A_k$ and
$A_{k+1}=A_k+J_k$. In the long run, what fraction of the pairs are juveniles? Does the answer depend on how many juveniles
and adults you start with? Example~\ref{exa:leslie} answers, and the golden ratio appears twice.
\end{think}

\section{Positive matrices: Perron's theorem}

We write $A\ge0$ when all entries of $A$ (a matrix or a vector) are non-negative, $A>0$ when they are all positive, and
$|x|$ for the vector of moduli of the entries of $x$. The \emph{spectral radius} of a square matrix is $\rho(A)=\max|\lambda|$
over its complex eigenvalues. Two remarks are used constantly: if $A\ge0$, then $|Ax|\le A|x|$ (triangle inequality in each
coordinate); and if $A>0$, $x\ge0$ and $x\ne0$, then $Ax>0$.

For $A\ge0$ and $x\ge0$, $x\ne0$, the \emph{Collatz--Wielandt function} is
\[
r(x)=\min_{i:\ x_i>0}\frac{(Ax)_i}{x_i},
\]
the largest $t$ with $Ax\ge tx$.

\begin{theorem}[Perron, 1907]\label{thm:perron}
Let $A>0$ be an $n\times n$ matrix and $\rho=\rho(A)$. Then:
\begin{enumerate}[label=(\alph*)]
\item $\rho>0$ is an eigenvalue of $A$, with an eigenvector $z>0$ (the \emph{Perron vector});
\item $\rho$ is an algebraically simple eigenvalue;
\item every other eigenvalue satisfies $|\lambda|<\rho$;
\item every eigenvector $x\ge0$ of $A$ is a multiple of $z$.
\end{enumerate}
Moreover $\rho=\max_{x\ge0,x\ne0}r(x)$ (Collatz--Wielandt formula).
\end{theorem}
\begin{proof}
\emph{(a) Existence.} Let $\Delta=\{x\ge0:\sum x_i=1\}$ and $r^*=\sup_\Delta r$ (finite: $r(x)\le\max_i\sum_ja_{ij}$). If
$Ax\ge tx$ then $A(Ax)\ge tAx$, so $r(Ax)\ge r(x)$: applying $A$ never decreases $r$. The set $K=\{Ax/\sum_i(Ax)_i:x\in\Delta\}$
is compact and consists of positive vectors, on which $r(y)=\min_i(Ay)_i/y_i$ is continuous. Since $r$ is unchanged by
scaling and $r(Ax)\ge r(x)$, $\sup_\Delta r=\sup_Kr$, which is attained at some $z\in K$, so $z>0$ and $Az\ge r^*z$. If
$Az\ne r^*z$, then $Az-r^*z\ge0$ is non-zero, so $A(Az-r^*z)>0$, that is $A(Az)>r^*Az$ strictly, and $r(Az)>r^*$, a
contradiction. Hence $Az=r^*z$ with $z>0$, and $r^*>0$ since $Az>0$.

\emph{$r^*=\rho$.} If $Ay=\lambda y$ with $y\ne0$, then $|\lambda||y|=|Ay|\le A|y|$, so $r(|y|)\ge|\lambda|$ and $|\lambda|\le
r^*$. As $r^*$ is itself an eigenvalue, $r^*=\rho$, which proves the Collatz--Wielandt formula.

\emph{(c)} Let $|\lambda|=\rho$ and $Ay=\lambda y$. Then $A|y|\ge\rho|y|$, and the equality argument of (a) applied to
$|y|$ shows $A|y|=\rho|y|$ and $|y|=A|y|/\rho>0$. So for each $i$, $\bigl|\sum_ja_{ij}y_j\bigr|=\sum_ja_{ij}|y_j|$ with all
$a_{ij}>0$: equality in the triangle inequality forces all $y_j$ to have the same argument, $y=e^{i\theta}|y|$. Then $\lambda
y=Ay=e^{i\theta}A|y|=\rho y$, so $\lambda=\rho$.

\emph{Geometric simplicity.} Let $y$ be an eigenvector for $\rho$; taking real and imaginary parts we may assume $y$ real,
and replacing $y$ by $-y$ that some $y_i>0$. Let $t=\min\{z_i/y_i:y_i>0\}$ and $w=z-ty$. Then $w\ge0$, $w$ has a zero
coordinate, and $Aw=\rho w$. If $w\ne0$, then $Aw>0$ while $\rho w$ has a zero coordinate: impossible. So $w=0$ and $y$ is
a multiple of $z$.

\emph{(b)} The matrix $A^{\mathsf T}>0$ has the same spectrum, hence by (a) a vector $w>0$ with $w^{\mathsf T}A=\rho w^{\mathsf
T}$. If $\rho$ were not algebraically simple, then, being geometrically simple, it would carry a Jordan block of size at least
$2$: there would be $u$ with $(A-\rho I)u=z$. Multiplying on the left by $w^{\mathsf T}$ gives $0=w^{\mathsf T}z>0$, a
contradiction.

\emph{(d)} If $Ax=\lambda x$ with $x\ge0$, $x\ne0$, then $\lambda w^{\mathsf T}x=w^{\mathsf T}Ax=\rho w^{\mathsf T}x$ and $w^{\mathsf
T}x>0$, so $\lambda=\rho$ and $x$ is a multiple of $z$ by geometric simplicity.
\end{proof}

The vector $w>0$ with $w^{\mathsf T}A=\rho w^{\mathsf T}$ is the \emph{left Perron vector}; we normalise $w^{\mathsf T}z=1$.

\section{Contraction and the power method}

The second proof explains \emph{why} the Perron vector wins: a positive matrix squeezes all positive directions towards one.
For positive vectors $x,y$, the ratios $x_i/y_i$ lie in an interval $[m(x/y),M(x/y)]$; we call its length the
\emph{oscillation} of $x/y$. It is zero exactly when $x$ is a multiple of $y$.

\begin{lemma}\label{lem:oscillation}
Let $A>0$, and let $c=\min a_{ij}/a_{kj}$ over all $i,j,k$, so $0<c\le1$. For positive $x,y$,
\[
\begin{gathered}
[m(Ax/Ay),M(Ax/Ay)]\subseteq[m(x/y),M(x/y)],\\
M(Ax/Ay)-m(Ax/Ay)\le(1-c^2)\bigl(M(x/y)-m(x/y)\bigr).
\end{gathered}
\]
\end{lemma}
\begin{proof}
Write $m=m(x/y)$ and $M=M(x/y)$, so that $p=A(x-my)$ and $q=A(My-x)$ are $\ge0$ with $p+q=(M-m)Ay$. The inclusion follows from
$(Ax)_i=m(Ay)_i+p_i\ge m(Ay)_i$ and similarly with $q$. For any $u\ge0$ and indices $i,k$, $(Au)_i=\sum_ja_{ij}u_j\ge c\sum_j
a_{kj}u_j=c(Au)_k$; in particular $(Ay)_i\le c^{-1}(Ay)_k$. So $\beta_i=q_i/(Ay)_i$ satisfies $\beta_i\ge c^2\beta_k$ for all
$i,k$. Since $(Ax)_i/(Ay)_i=M-\beta_i$, the oscillation of $Ax/Ay$ is $\max\beta-\min\beta\le(1-c^2)\max\beta\le(1-c^2)(M-m)$,
as $\beta_i\le M-m$ (because $p_i\ge0$).
\end{proof}

\begin{theorem}[convergence of the power method]\label{thm:powerconv}
Let $A>0$ with Perron vectors $z,w$, $w^{\mathsf T}z=1$. For every $x\ge0$, $x\ne0$, $A^kx/\rho^k\to(w^{\mathsf T}x)\,z$, and
$A^k/\rho^k\to zw^{\mathsf T}$. The convergence is geometric.
\end{theorem}
\begin{proof}
Replacing $x$ by $Ax$ we may assume $x>0$. Apply Lemma~\ref{lem:oscillation} repeatedly with $y=z$, noting $A^kz=\rho^kz$:
the ratios $(A^kx)_i/(\rho^kz_i)$ lie in nested intervals whose lengths are at most $(1-c^2)^k$ times the first. So they
converge, uniformly in $i$ and geometrically, to a common limit $\gamma$: $A^kx/\rho^k\to\gamma z$. Since $w^{\mathsf T}A^kx/
\rho^k=w^{\mathsf T}x$ for all $k$, $\gamma=w^{\mathsf T}x$. Applying this to the basis vectors gives the matrix statement.
\end{proof}

So for any positive starting vector, $A^kx/\|A^kx\|$ converges to the Perron vector: this is the \emph{power method}, and
PageRank is computed this way. The proof shows a rate $1-c^2$; the true rate is $|\lambda_2|/\rho$, where $\lambda_2$ is the
second largest eigenvalue in modulus, as one sees from the Jordan form.

\begin{remark}
Lemma~\ref{lem:oscillation} is a version of a beautiful fact found by Garrett Birkhoff (1957) and Hans Samelson: in
\emph{Hilbert's projective metric} $d(x,y)=\log\frac{M(x/y)}{m(x/y)}$ on the positive cone, every positive matrix is a strict
contraction, with the explicit contraction factor $\tanh(\Delta/4)$, where $\Delta$ is the diameter of the image of the cone.
Perron's theorem then becomes the Banach fixed point theorem; this point of view extends to infinite dimension and to
non-linear maps.
\end{remark}

\section{Non-negative matrices: Frobenius's theorem}

Zeros change the picture. The matrix $\bigl(\begin{smallmatrix}0&1\\1&0\end{smallmatrix}\bigr)$ has eigenvalues $\pm1$, both
of modulus $\rho$; the matrix $\bigl(\begin{smallmatrix}1&1\\0&1\end{smallmatrix}\bigr)$ has a Jordan block at $\rho$, and its
eigenvector $(1,0)$ is not positive. The right notion is connectivity. To $A\ge0$ of size $n$ associate the directed graph
$\Gamma(A)$ on vertices $1,\dots,n$ with an edge $i\to j$ when $a_{ij}>0$. Then $(A^k)_{ij}>0$ if and only if there is a walk of
length $k$ from $i$ to $j$, since $(A^k)_{ij}$ is a sum of products $a_{ii_1}a_{i_1i_2}\cdots a_{i_{k-1}j}$ over such walks.

\begin{definition}\label{def:irreducible}
$A\ge0$ is \emph{irreducible} if $\Gamma(A)$ is strongly connected: there is a walk from every vertex to every vertex. Its
\emph{period} $h$ is the greatest common divisor of the lengths of the closed walks of $\Gamma(A)$. It is \emph{primitive} if
$A^k>0$ for some $k$.
\end{definition}

Since a path between distinct vertices visits at most $n$ vertices, $A$ is irreducible if and only if $(I+A)^{n-1}>0$: the
$(i,j)$ entry of $(I+A)^{n-1}=\sum_k\binom{n-1}kA^k$ is positive exactly when some walk of length at most $n-1$ joins $i$ to $j$.

\begin{theorem}[Frobenius, 1912]\label{thm:frobenius}
Let $A\ge0$ be irreducible, of period $h$, and $\rho=\rho(A)$. Then $\rho>0$ (if $n\ge2$) is an algebraically simple
eigenvalue with an eigenvector $z>0$, and every eigenvector $x\ge0$ is a multiple of $z$. The eigenvalues of modulus $\rho$ are
exactly $\rho\,\omega^k$, $k=0,\dots,h-1$, where $\omega=e^{2\pi i/h}$; each is simple, and the whole spectrum is invariant under
multiplication by $\omega$.
\end{theorem}
\begin{proof}
\emph{Perron part.} $B=(I+A)^{n-1}>0$ commutes with $A$. By Theorem~\ref{thm:perron}, $B$ has a one-dimensional eigenspace for
$\rho(B)$, spanned by $z>0$; as $A$ commutes with $B$, it preserves this line, so $Az=rz$ for some $r$, and $r\ge0$ since $Az\ge0$.
The same argument with $A^{\mathsf T}$ gives $w>0$ and $r'$ with $w^{\mathsf T}A=r'w^{\mathsf T}$, and $r'w^{\mathsf T}z=w^{\mathsf
T}Az=rw^{\mathsf T}z$ gives $r'=r$. If $Ay=\lambda y$, then $|\lambda||y|\le A|y|$, and multiplying by $w^{\mathsf T}$,
$|\lambda|w^{\mathsf T}|y|\le rw^{\mathsf T}|y|$, so $|\lambda|\le r$: hence $r=\rho$. If $\rho=0$, then $Az=0$ with $z>0$
forces $A=0$, which is not irreducible for $n\ge2$. The eigenvalues of $B$ are the $(1+\lambda)^{n-1}$, with multiplicities,
and $|1+\lambda|\le1+|\lambda|\le1+\rho$ with equality only for $\lambda=\rho$; since $(1+\rho)^{n-1}=\rho(B)$ is simple for $B$,
$\rho$ is simple for $A$. Non-negative eigenvectors are multiples of $z$ by the argument of Theorem~\ref{thm:perron}(d).

\emph{The peripheral spectrum.} Let $Ay=\lambda y$ with $\lambda=\rho e^{i\varphi}$. As above, $A|y|\ge\rho|y|$, and
$w^{\mathsf T}(A|y|-\rho|y|)=0$ with $w>0$ forces $A|y|=\rho|y|$, so $|y|$ is a multiple of $z$, positive. Write $y_j=e^{i\theta_j}
|y_j|$. Then for each $i$, $\sum_ja_{ij}|y_j|e^{i\theta_j}=\rho e^{i\varphi}e^{i\theta_i}|y_i|=e^{i(\varphi+\theta_i)}\sum_ja_{ij}|y_j|$;
equality in the triangle inequality gives
\[
e^{i\theta_j}=e^{i\varphi}e^{i\theta_i}\qquad\text{whenever }a_{ij}>0.
\]
Following a closed walk of length $\ell$ multiplies $e^{i\theta}$ by $e^{i\ell\varphi}$ and returns to the start, so $e^{i\ell
\varphi}=1$ for every closed walk length $\ell$, hence $e^{ih\varphi}=1$ (the period is an integer combination of such lengths).
So the peripheral eigenvalues are among the $\rho\omega^k$.

Conversely, Proposition~\ref{prop:levels} below gives a function $\ell:\{1,\dots,n\}\to\Z/h$ with $\ell(j)=\ell(i)+1$ for every
edge $i\to j$. With $D=\operatorname{diag}(\omega^{\ell(1)},\dots,\omega^{\ell(n)})$, $(D^{-1}AD)_{ij}=\omega^{\ell(j)-\ell(i)}a_{ij}=
\omega a_{ij}$. So $A$ is similar to $\omega A$: the spectrum, with multiplicities, is invariant under multiplication by $\omega$,
and each $\rho\omega^k$ is an eigenvalue as simple as $\rho$.
\end{proof}

\begin{proposition}\label{prop:levels}
Let $\Gamma$ be strongly connected with period $h$. Fix a vertex $v_0$. The length of a walk from $v_0$ to $j$ is
well defined modulo $h$; calling it $\ell(j)$, every edge $i\to j$ satisfies $\ell(j)=\ell(i)+1$ in $\Z/h$. So the vertices
fall into $h$ classes, and every edge goes from one class to the next, cyclically.
\end{proposition}
\begin{proof}
Let two walks from $v_0$ to $j$ have lengths $p$ and $q$, and let a walk from $j$ back to $v_0$ have length $s$. Then $p+s$ and
$q+s$ are closed walk lengths, so $h$ divides both and $p\equiv q\pmod h$. A walk to $i$ followed by the edge $i\to j$ is a walk
to $j$, which gives $\ell(j)=\ell(i)+1$.
\end{proof}

\begin{figure}[ht]
\centering
\fig{ch07-period}
\caption{A strongly connected digraph of period $3$. Every edge goes from one colour class to the next, so all closed walks
have lengths divisible by $3$; the adjacency matrix has three eigenvalues of maximal modulus, $\rho$, $\rho\omega$, $\rho\omega^2$.}
\label{fig:period}
\end{figure}

\section{Primitive matrices, Wielandt's bound}

\begin{theorem}\label{thm:primitive}
For irreducible $A\ge0$, the following are equivalent: (i) $A$ is primitive; (ii) the period is $h=1$; (iii) $\rho$ is the
only eigenvalue of modulus $\rho$. For a primitive $A$, $A^k/\rho^k\to zw^{\mathsf T}$, where $z,w$ are the right and left Perron
vectors with $w^{\mathsf T}z=1$.
\end{theorem}
\begin{proof}
(ii)$\Leftrightarrow$(iii) is Theorem~\ref{thm:frobenius}. (i)$\Rightarrow$(iii): if $A^k>0$, then $\rho^k$ is a simple eigenvalue
of $A^k$ and the only one of modulus $\rho^k$ (Theorem~\ref{thm:perron}). If $\lambda\ne\rho$ had $|\lambda|=\rho$, then $\lambda^k$
would have modulus $\rho^k$, hence equal $\rho^k$, which would then have multiplicity at least $2$ in $A^k$.

(ii)$\Rightarrow$(i): fix a vertex $i$ and let $S$ be the set of lengths of closed walks at $i$. It is closed under addition. Its
gcd is $1$: for any closed walk of length $\ell$ through a vertex $j$, going from $i$ to $j$ (length $a$), around (length $\ell$)
and back (length $b$) shows $a+b\in S$ and $a+b+\ell\in S$, so $\gcd S$ divides every closed walk length, hence divides $h=1$.
A set of positive integers closed under addition with gcd $1$ contains all large integers: by Bezout, $1=\sum_tc_ts_t$ with
$s_t\in S$ and $c_t\in\Z$; letting $P$ be the sum of the positive terms and $N$ minus the sum of the negative ones, both in $S$
(or zero) with $P-N=1$; then for $m\ge N^2$, write $m=qN+r$ with $0\le r<N$ and $q\ge N$, so that $m=rP+(q-r)N\in S$. Hence
$(A^k)_{ii}>0$ for all large $k$, for every $i$, and then $(A^k)_{ij}>0$ for all large $k$ by appending a fixed path from $i$ to $j$.

For the limit, $B=A^N>0$ for some $N$, with the same Perron vectors and $\rho(B)=\rho^N$. For $k=mN+s$ with $0\le s<N$,
Theorem~\ref{thm:powerconv} for $B$ gives $A^k/\rho^k=(B^m/\rho^{mN})A^s/\rho^s\to zw^{\mathsf T}A^s/\rho^s=zw^{\mathsf T}$.
\end{proof}

How large must $k$ be? Wielandt stated the sharp answer in 1950 without proof; the first published proofs came in 1957
and 1958 (Rosenblatt, Holladay and Varga), and Wielandt's own unpublished proof appeared in print only in 2002.

\begin{theorem}[Wielandt]\label{thm:wielandt}
If $A\ge0$ of size $n\ge2$ is primitive, then $A^{(n-1)^2+1}>0$.
\end{theorem}
\begin{proof}
Let $s$ be the length of a shortest cycle $C$ of $\Gamma=\Gamma(A)$. Then $s\le n-1$: otherwise every cycle would have length $n$,
and the period would be $n\ge2$ (every closed walk decomposes into cycles). The matrix $B=A^s$ is primitive (a power of a
primitive matrix is primitive, since $A$ has no zero row and $A^m>0$ for all large $m$), so $\Gamma(B)$ is strongly connected, and
every vertex $c$ of $C$ has a loop in $\Gamma(B)$ (go once around $C$). From $c$, every vertex $j$ is reached in $\Gamma(B)$ by a
path of length at most $n-1$, which the loop at $c$ pads to length exactly $n-1$: $(B^{n-1})_{cj}=(A^{s(n-1)})_{cj}>0$.

Now let $i$ be any vertex. A shortest path from $i$ to $C$ has length $t\le n-s$ (it visits distinct vertices outside $C$
before reaching $C$); continue along $C$ for $n-s-t$ more steps, to get a walk of length exactly $n-s$ from $i$ to some vertex
$c$ of $C$. Then $(A^{n-s+s(n-1)})_{ij}\ge(A^{n-s})_{ic}(A^{s(n-1)})_{cj}>0$ for every $j$. Since $n-s+s(n-1)=n+s(n-2)\le
n+(n-1)(n-2)=(n-1)^2+1$, and $A^{k}>0$ implies $A^{k+1}=A\cdot A^k>0$ (no zero row), the theorem follows.
\end{proof}

The bound is attained by the \emph{Wielandt matrix}, the graph $1\to2\to\dots\to n\to1$ with one extra edge $n\to2$; the notebook
checks this for $n\le10$.

\section{Markov chains and PageRank}

A \emph{stochastic} matrix $P\ge0$ has rows summing to $1$: $P\mathbf 1=\mathbf 1$. It describes a Markov chain, $p_{ij}$ being the
probability to move from state $i$ to state $j$, and a probability distribution $\mu$ (a row vector) evolves as $\mu\mapsto\mu
P$. The vector $\mathbf 1>0$ is a non-negative eigenvector for the eigenvalue $1$, so $\rho(P)=1$ when $P$ is irreducible
(Theorem~\ref{thm:frobenius}: non-negative eigenvectors belong to $\rho$), and indeed for every stochastic $P$, since $\|Px\|_\infty
\le\|x\|_\infty$.

\begin{corollary}[ergodic theorem for Markov chains]\label{cor:markov}
Let $P$ be a stochastic matrix that is irreducible with period $1$. There is a unique probability vector $\pi$ with $\pi P=\pi$;
it is positive, and for every initial distribution $\mu$, $\mu P^k\to\pi$. More precisely $P^k\to\mathbf 1\pi$, the matrix all of
whose rows are $\pi$.
\end{corollary}
\begin{proof}
Here $\rho=1$, $z=\mathbf 1$, and the left Perron vector normalised by $w^{\mathsf T}\mathbf 1=1$ is $\pi$, unique and positive by
Theorem~\ref{thm:frobenius}. Theorem~\ref{thm:primitive} gives $P^k\to\mathbf 1\pi$, and $\mu P^k\to(\mu\mathbf 1)\pi=\pi$.
\end{proof}

Tome~3, Chapter~8 proves the same theorem by coupling, with probabilistic intuition and explicit rates; here it is a statement
about eigenvalues. The rate of convergence is governed by $|\lambda_2|$, the \emph{spectral gap} $1-|\lambda_2|$ measuring how fast
the chain forgets where it started.

\begin{example}[PageRank]\label{exa:pagerank}
Model a surfer on the web: from page $i$ with $d_i$ outgoing links, follow one at random ($p_{ij}=1/d_i$ for each link
$i\to j$; a page with no links is given links to all pages). The importance of a page should be the long-run fraction of
time the surfer spends there, the stationary distribution of $P$. But the web graph is not strongly connected, so Brin and
Page added teleportation: with probability $1-\alpha$, the surfer jumps to a uniformly random page. The \emph{Google matrix}
\[
G=\alpha P+\frac{1-\alpha}n\mathbf 1\mathbf 1^{\mathsf T},\qquad\alpha=0.85,
\]
is stochastic and positive, so Perron's theorem gives a unique positive PageRank vector $\pi=\pi G$, and the power method
converges at rate at most $\alpha$ (Exercise~\ref{ex:pagerankrate}), independently of the size of the web: about $50$
iterations give six digits. Figure~\ref{fig:pagerank} shows a small example. Page $C$, with six incoming links, ranks first. Page $A$
has only two, but one of them comes from $C$, which has no other outgoing link and passes all its weight to $A$: so $A$ ranks
almost as high as $C$. Page $B$ also has two incoming links, one from the important page $A$, and ranks third. Being linked by
important pages matters more than being linked often.
\end{example}

\begin{figure}[ht]
\centering
\fig{ch07-pagerank}
\caption{PageRank of an eight-page web with $\alpha=0.85$; the area of each node is proportional to its rank. The values
are $C:0.351$, $A:0.329$, $B:0.169$, and between $0.019$ and $0.040$ for the others.}
\label{fig:pagerank}
\end{figure}

\section{Populations}

\begin{example}[Leslie models]\label{exa:leslie}
Divide a female population into age classes $1,\dots,n$. Each step, a female of class $i$ gives birth on average to $f_i\ge0$
daughters, and survives into class $i+1$ with probability $s_i\in(0,1]$. The vector $x_k$ of class sizes evolves by
$x_{k+1}=Lx_k$, with the \emph{Leslie matrix}
\[
L=\begin{pmatrix}f_1&f_2&\cdots&f_{n-1}&f_n\\s_1&0&\cdots&0&0\\0&s_2&&&0\\&&\ddots&&\vdots\\0&0&\cdots&s_{n-1}&0\end{pmatrix}.
\]
If $f_n>0$, the graph $\Gamma(L)$ contains the cycle $1\to n\to n-1\to\dots\to2\to1$ (the edge $1\to j$ comes from
$f_j>0$ and the edge $j+1\to j$ from $s_j>0$), so $L$ is irreducible; if moreover two consecutive $f_i,f_{i+1}$ are positive,
there are cycles $1\to i\to\dots\to1$ of lengths $i$ and $i+1$, the period is $1$ and $L$ is primitive. Then $L^kx_0/\rho^k\to(w^{\mathsf T}x_0)z$: whatever the initial population,
it eventually grows by the factor $\rho$ each step, with the \emph{stable age distribution} $z$. The characteristic equation
is the \emph{Euler--Lotka equation} $\sum_if_i\ell_i\rho^{-i}=1$, with $\ell_i=s_1\cdots s_{i-1}$ the probability to reach class
$i$ (Exercise~\ref{ex:leslie}).

Fibonacci's rabbits, counted in pairs, are a Leslie model with two classes, juveniles and adults, $f_1=0$,
$f_2=1$, $s_1=1$, and one extra entry $1$ in position $(2,2)$ because adults remain adults:
$L=\bigl(\begin{smallmatrix}0&1\\1&1\end{smallmatrix}\bigr)$ in the coordinates $(J,A)$ of the opening question, with
$\rho=\varphi=\frac{1+\sqrt5}2$ and Perron vector $z=(1,\varphi)$ (indeed $Lz=(\varphi,1+\varphi)=\varphi z$). This $L$ is
primitive ($L^2>0$), so the fraction of juveniles tends to $\frac1{1+\varphi}=\frac1{\varphi^2}\approx0.382$, whatever the
start: the golden ratio twice, as the growth factor and in the stable proportions.
\end{example}

\begin{beauty}
A positive matrix knows which way is up. However you start, as long as you start inside the positive cone, repeated
multiplication drags every direction to the same one, the Perron vector, and this single geometric fact predicts the age
pyramid of a rat colony, the rank of a web page, and the long-run frequencies of a card shuffle. When zeros break the
positivity, the leading eigenvalues arrange themselves as a regular polygon whose number of sides is the gcd of the cycle
lengths: arithmetic of the graph becomes geometry of the spectrum.
\end{beauty}

\section{Exercises}

\begin{exercise}[$\star$]\label{ex:perron2}
Let $A=\bigl(\begin{smallmatrix}a&b\\c&d\end{smallmatrix}\bigr)>0$. Compute $\rho$ and the Perron vector, and check directly
that the other eigenvalue has smaller modulus.
\end{exercise}

\begin{exercise}[$\star$]\label{ex:rowsum}
For $A\ge0$ irreducible, show that $\min_i\sum_ja_{ij}\le\rho\le\max_i\sum_ja_{ij}$, with equality on either side only if all row
sums are equal.
\end{exercise}

\begin{exercise}[$\star\star$]\label{ex:minmaxcw}
For $A>0$, show the dual Collatz--Wielandt formula $\rho=\min_{x>0}\max_i\frac{(Ax)_i}{x_i}$, and deduce that any $x>0$ gives
$\min_i\frac{(Ax)_i}{x_i}\le\rho\le\max_i\frac{(Ax)_i}{x_i}$.
\end{exercise}

\begin{exercise}[$\star\star$]\label{ex:monotone}
Let $0\le A\le B$ entrywise, with $B$ irreducible. Show that $\rho(A)\le\rho(B)$, with equality only if $A=B$.
\end{exercise}

\begin{exercise}[$\star\star$]\label{ex:nonneg}
Let $A\ge0$ be arbitrary. Show that $\rho(A)$ is an eigenvalue of $A$ with an eigenvector $x\ge0$, $x\ne0$. (Apply
Exercise~\ref{ex:monotone} to $A+\varepsilon J$, where $J$ is the all-ones matrix, and let $\varepsilon\to0$.)
\end{exercise}

\begin{exercise}[$\star$]\label{ex:cycleperiod}
Find the period and the eigenvalues of the adjacency matrix of the directed cycle on $n$ vertices, and of the adjacency matrix
of a connected bipartite graph (with each edge in both directions). Show that the spectrum of a connected bipartite graph is
symmetric about $0$.
\end{exercise}

\begin{exercise}[$\star$]\label{ex:diagpos}
Show that an irreducible $A\ge0$ with a positive diagonal entry is primitive, and that if all diagonal entries are positive then
$A^{n-1}>0$.
\end{exercise}

\begin{exercise}[$\star$]\label{ex:twostate}
For the two-state chain $P=\bigl(\begin{smallmatrix}1-a&a\\b&1-b\end{smallmatrix}\bigr)$ with $0<a,b<1$, find $\pi$ and an
explicit formula for $P^k$. What is the rate of convergence?
\end{exercise}

\begin{exercise}[$\star\star$]\label{ex:pagerankrate}
Show that every eigenvalue $\lambda\ne1$ of the Google matrix $G=\alpha P+(1-\alpha)\mathbf 1v^{\mathsf T}$ ($P$ stochastic, $v$ a
probability vector) satisfies $|\lambda|\le\alpha$.
\end{exercise}

\begin{exercise}[$\star\star$]\label{ex:leslie}
Show that the characteristic polynomial of the Leslie matrix is $X^n-\sum_if_i\ell_iX^{n-i}$, that it has exactly one
positive root when some $f_i>0$, and that this root is $\rho(L)$. Compute the Perron vector, and show that $\rho>1$ if and only
if the \emph{net reproductive rate} $\sum_if_i\ell_i$ exceeds $1$.
\end{exercise}

\begin{exercise}[$\star\star$]\label{ex:doublystoch}
Let $P$ be a doubly stochastic primitive matrix. Show that $P^k\to\frac1nJ$. Deduce that repeating any fixed card shuffle
whose transition matrix is primitive makes all $n!$ orders of the deck equally likely in the limit.
\end{exercise}

\begin{exercise}[$\star\star\star$]\label{ex:wielandtsharp}
Let $W$ be the Wielandt matrix (edges $i\to i+1$ for $i<n$, $n\to1$ and $n\to2$). Show that every closed walk at vertex $1$
has length $an+b(n-1)$ with $a\ge1$ and $b\ge0$, and deduce that $(W^{(n-1)^2})_{11}=0$, so that Theorem~\ref{thm:wielandt} is
sharp.
\end{exercise}

\section{Hints}

\begin{hint}{ex:perron2}
$\rho=\frac{a+d+\sqrt{(a-d)^2+4bc}}2$.
\end{hint}
\begin{hint}{ex:rowsum}
Look at the row of $Az=\rho z$ where $z_i$ is largest, and the one where it is smallest.
\end{hint}
\begin{hint}{ex:minmaxcw}
For $x>0$, apply the left Perron vector to $Ax\le Mx$.
\end{hint}
\begin{hint}{ex:monotone}
If $Ay=\lambda y$ then $|\lambda||y|\le B|y|$; multiply by the left Perron vector of $B$. For equality, look at what is forced.
\end{hint}
\begin{hint}{ex:nonneg}
$\rho(A+\varepsilon J)$ decreases as $\varepsilon\downarrow0$; take a limit of normalised Perron vectors.
\end{hint}
\begin{hint}{ex:cycleperiod}
For a bipartite graph, the classes of Proposition~\ref{prop:levels} are the two sides.
\end{hint}
\begin{hint}{ex:diagpos}
A loop at a vertex makes $1$ a closed walk length. With all loops, walks can be padded to any longer length.
\end{hint}
\begin{hint}{ex:twostate}
The eigenvalues are $1$ and $1-a-b$.
\end{hint}
\begin{hint}{ex:pagerankrate}
A left eigenvector $x$ for $\lambda\ne1$ satisfies $x^{\mathsf T}\mathbf 1=0$.
\end{hint}
\begin{hint}{ex:leslie}
Expand along the first row, or solve $Lz=\lambda z$ from the bottom up. Descartes's rule, or monotonicity of $\sum f_i\ell_i
\lambda^{-i}$.
\end{hint}
\begin{hint}{ex:doublystoch}
$\mathbf 1$ is both a left and a right Perron vector.
\end{hint}
\begin{hint}{ex:wielandtsharp}
Vertex $n$ is the only vertex with two outgoing edges. Reduce $(n-1)^2=an+b(n-1)$ modulo $n-1$.
\end{hint}

\section{Solutions}

\begin{solution}{ex:perron2}
\idea Solve the quadratic and compare the two roots.
\writeup $\chi_A=X^2-(a+d)X+ad-bc$, with discriminant $(a-d)^2+4bc>0$, so the roots are real and distinct,
$\lambda_\pm=\frac{a+d\pm\sqrt{(a-d)^2+4bc}}2$. Since $\sqrt{(a-d)^2+4bc}>|a-d|$, we get $\lambda_+>\max(a,d)>0$, and
$|\lambda_-|<\lambda_+$ because $\lambda_++\lambda_-=a+d>0$. An eigenvector for $\lambda_+$ is $(b,\lambda_+-a)$, positive since
$\lambda_+>a$.
\end{solution}

\begin{solution}{ex:rowsum}
\idea Read the equation $Az=\rho z$ in the coordinate where $z$ is largest: $\rho z_i$ is an average-like combination of
smaller values.
\writeup Let $z>0$ be the Perron vector and $z_i=\max z$. Then $\rho z_i=\sum_ja_{ij}z_j\le z_i\sum_ja_{ij}\le z_i\max_k\sum_j
a_{kj}$. At the row where $z$ is smallest, the same argument gives the lower bound. Equality $\rho=\max$ row sum forces, in the
row $i$, $a_{ij}z_j=a_{ij}z_i$ for all $j$, i.e.\ $z_j=z_i$ for the neighbours $j$ of $i$, and the row sum of $i$ to be maximal.
The neighbours then also achieve the maximum of $z$, and the argument propagates along edges; by irreducibility $z$ is constant,
and then $A\mathbf 1=\rho\mathbf 1$: all row sums equal $\rho$.
\end{solution}

\begin{solution}{ex:minmaxcw}
\idea Pair any test vector with the left Perron vector: $w^{\mathsf T}Ax=\rho w^{\mathsf T}x$ is the bridge.
\writeup Let $x>0$ and $M=\max_i(Ax)_i/x_i$, so $Ax\le Mx$. Then $\rho w^{\mathsf T}x=w^{\mathsf T}Ax\le Mw^{\mathsf T}x$, and
$w^{\mathsf T}x>0$, so $\rho\le M$. Equality for $x=z$. The lower bound is the Collatz--Wielandt formula of
Theorem~\ref{thm:perron}, or the same argument with $Ax\ge mx$.
\end{solution}

\begin{solution}{ex:monotone}
\idea Use $B$'s left Perron vector as a weighting that turns vector inequalities into number inequalities.
\writeup If $Ay=\lambda y$, $y\ne0$, then $|\lambda||y|\le A|y|\le B|y|$. With $w>0$ the left Perron vector of $B$,
$|\lambda|w^{\mathsf T}|y|\le w^{\mathsf T}B|y|=\rho(B)w^{\mathsf T}|y|$, so $|\lambda|\le\rho(B)$. If $\rho(A)=\rho(B)$, take
$\lambda$ with $|\lambda|=\rho(B)$: all inequalities are equalities, $w^{\mathsf T}(B|y|-\rho(B)|y|)=0$ gives $B|y|=\rho(B)|y|$, so
$|y|=z>0$ (Frobenius), and $w^{\mathsf T}(B-A)z=0$ with $w,z>0$ and $B-A\ge0$ forces $A=B$.
\end{solution}

\begin{solution}{ex:nonneg}
\idea Positive perturbations have Perron vectors; compactness of the simplex lets us pass to the limit, and monotonicity
controls the eigenvalue.
\writeup $A_\varepsilon=A+\varepsilon J>0$ has Perron root $\rho_\varepsilon$ and Perron vector $z_\varepsilon$ with $\sum_i
(z_\varepsilon)_i=1$. By Exercise~\ref{ex:monotone}, $\rho_\varepsilon$ decreases as $\varepsilon$ decreases, so it has a limit $r$
as $\varepsilon\to0$, and the same argument (with $B=A_\varepsilon$, which is irreducible) gives $\rho(A)\le\rho_\varepsilon$, so
$\rho(A)\le r$. By compactness, some sequence $z_{\varepsilon_k}$ converges to $x\ge0$ with $\sum x_i=1$, and passing to the limit in
$A_\varepsilon z_\varepsilon=\rho_\varepsilon z_\varepsilon$ gives $Ax=rx$. So $r$ is an eigenvalue of $A$, $r\le\rho(A)$, and
therefore $r=\rho(A)$ with eigenvector $x\ge0$.
\end{solution}

\begin{solution}{ex:cycleperiod}
\idea A directed $n$-cycle is the permutation matrix of a cyclic shift; for bipartite graphs, use the similarity of
Theorem~\ref{thm:frobenius} with $h=2$.
\writeup The cycle has all closed walks of length divisible by $n$, so $h=n$; its matrix $C$ satisfies $C^n=I$ and its
eigenvalues are the $n$-th roots of unity, each once (eigenvectors $(\omega^{jk})_j$): exactly the peripheral spectrum predicted,
with $\rho=1$. A connected bipartite graph has all closed walks of even length and contains closed walks of length $2$ (an
edge back and forth), so $h=2$; Theorem~\ref{thm:frobenius} gives that $A$ is similar to $-A$ (explicitly, $D=\operatorname{diag}
(\pm1)$ according to the side), so the spectrum is symmetric about $0$.
\end{solution}

\begin{solution}{ex:diagpos}
\idea A loop is a closed walk of length $1$, so the gcd is $1$; loops at every vertex let any short walk wait.
\writeup If $a_{ii}>0$, then $1$ is a closed walk length, so $h=1$ and $A$ is primitive by Theorem~\ref{thm:primitive}. If
all $a_{ii}>0$, take any $i,j$: by irreducibility there is a path from $i$ to $j$ of some length $k\le n-1$, and padding it with
$n-1-k$ loops at $j$ gives a walk of length exactly $n-1$. So $(A^{n-1})_{ij}>0$ for all $i,j$.
\end{solution}

\begin{solution}{ex:twostate}
\idea Two states, two eigenvalues: $1$ and the trace minus $1$.
\writeup $\pi=\bigl(\frac b{a+b},\frac a{a+b}\bigr)$ satisfies $\pi P=\pi$. The eigenvalues are $1$ and $\lambda=1-a-b\in(-1,1)$,
and
\[
P^k=\mathbf 1\pi+\lambda^k\frac1{a+b}\begin{pmatrix}a&-a\\-b&b\end{pmatrix},
\]
as one checks for $k=0$ and by noting that the second matrix $E$ satisfies $PE=EP=\lambda E$ and $\mathbf 1\pi E=0$. So
$P^k\to\mathbf 1\pi$ at rate $|1-a-b|^k$: fast when $a+b$ is near $1$, slow for a ``sticky'' chain ($a,b$ small) or an almost
periodic one ($a,b$ near $1$).
\end{solution}

\begin{solution}{ex:pagerankrate}
\idea Teleportation is a rank-one perturbation that acts trivially on vectors orthogonal to $\mathbf 1$.
\writeup Let $x^{\mathsf T}G=\lambda x^{\mathsf T}$ with $\lambda\ne1$, $x\ne0$. Since $G\mathbf 1=\mathbf 1$, $\lambda x^{\mathsf
T}\mathbf 1=x^{\mathsf T}G\mathbf 1=x^{\mathsf T}\mathbf 1$, so $x^{\mathsf T}\mathbf 1=0$. Then $x^{\mathsf T}G=\alpha x^{\mathsf T}P+
(1-\alpha)(x^{\mathsf T}\mathbf 1)v^{\mathsf T}=\alpha x^{\mathsf T}P$, so $\lambda/\alpha$ is an eigenvalue of $P$ (if $\alpha>0$),
and $|\lambda/\alpha|\le\rho(P)=1$.
\end{solution}

\begin{solution}{ex:leslie}
\idea An eigenvector of a Leslie matrix is forced by the survival rows; the first row then gives one scalar equation.
\writeup Suppose $Lz=\lambda z$ with $\lambda\ne0$. Rows $2,\dots,n$ give $z_{i+1}=s_iz_i/\lambda$, so $z_i=\ell_i\lambda^{-(i-1)}z_1$.
The first row gives $\lambda z_1=\sum_if_iz_i=z_1\sum_if_i\ell_i\lambda^{1-i}$, i.e.\ $\sum_if_i\ell_i\lambda^{-i}=1$, which is
$\lambda^n-\sum_if_i\ell_i\lambda^{n-i}=0$ after multiplying by $\lambda^n$. Both sides are monic of degree $n$ and the
eigenvalue equation must be the characteristic polynomial (expand $\det(XI-L)$ along the first row to confirm). For $\lambda>0$,
$g(\lambda)=\sum_if_i\ell_i\lambda^{-i}$ is strictly decreasing from $+\infty$ to $0$ (some $f_i\ell_i>0$), so there is
exactly one positive root $\lambda_0$, with eigenvector $z=(\ell_i\lambda_0^{1-i})_i>0$. A non-negative eigenvector belongs to
$\rho$ (Exercise~\ref{ex:nonneg} gives one for $\rho$, and its eigenvalue is a non-negative real root), so $\lambda_0=\rho$ when
$\rho>0$. Finally $\rho>1$ iff $g(1)>g(\rho)=1$, i.e.\ $\sum_if_i\ell_i>1$.
\end{solution}

\begin{solution}{ex:doublystoch}
\idea Doubly stochastic means $\mathbf 1$ is invariant on both sides, so both Perron vectors are uniform.
\writeup $P\mathbf 1=\mathbf 1$ and $\mathbf 1^{\mathsf T}P=\mathbf 1^{\mathsf T}$, so $z=\mathbf 1$ and $w=\frac1n\mathbf 1$ with
$w^{\mathsf T}z=1$, and Theorem~\ref{thm:primitive} gives $P^k\to\frac1n\mathbf 1\mathbf 1^{\mathsf T}$. A card shuffle is a
random permutation applied to the deck; its transition matrix on the $n!$ orders has entries $p_{\sigma\tau}=\Pr(\text{shuffle}=
\tau\sigma^{-1})$, and each column sums to $1$ as well as each row (for fixed $\tau$, $\sigma\mapsto\tau\sigma^{-1}$ is a
bijection). So it is doubly stochastic, and if primitive, the distribution of the deck tends to uniform.
\end{solution}

\begin{solution}{ex:wielandtsharp}
\idea Vertex $n$ is the only branching point, so walks are made of two loops; then it is a coin problem with coins $n$ and
$n-1$, and the condition $a\ge1$ (we must come back to $1$) is what makes $(n-1)^2$ unreachable.
\writeup Every vertex $i<n$ has the single outgoing edge $i\to i+1$. So a closed walk at $1$ goes $1\to2\to\dots\to n$ and then,
each time it is at $n$, either returns to $1$ (closing an excursion of length $n$ since the last visit to $1$) or goes around
$n\to2\to\dots\to n$ (length $n-1$). It ends at $1$, so it consists of $a\ge1$ excursions through the edge $n\to1$ and $b\ge0$ extra
loops through $n\to2$, and its length is $an+b(n-1)$. Suppose $(n-1)^2=an+b(n-1)$. Since $n\equiv1\pmod{n-1}$, reducing modulo
$n-1$ gives $a\equiv0$, so $a\ge n-1$ and $an\ge n(n-1)>(n-1)^2$, a contradiction. Hence there is no closed walk of length
$(n-1)^2$ at $1$, and $(W^{(n-1)^2})_{11}=0$. The notebook confirms that $W^{(n-1)^2+1}>0$.
\end{solution}

\section*{Sources and further reading}

\begin{itemize}
\item O.~Perron, \emph{Zur Theorie der Matrices}, Mathematische Annalen 64 (1907), 248--263; G.~Frobenius, \emph{\"Uber Matrizen aus
nicht negativen Elementen}, Sitzungsberichte der Preussischen Akademie (1912), 456--477.
\item L.~Collatz, \emph{Einschlie\ss ungssatz f\"ur die charakteristischen Zahlen von Matrizen}, Mathematische Zeitschrift 48
(1942), 221--226; H.~Wielandt, \emph{Unzerlegbare, nicht negative Matrizen}, Mathematische Zeitschrift 52 (1950), 642--648.
\item G.~Birkhoff, \emph{Extensions of Jentzsch's theorem}, Transactions of the AMS 85 (1957), 219--227 (Hilbert's metric).
\item H.~Schneider, \emph{Wielandt's proof of the exponent inequality for primitive nonnegative matrices}, Linear Algebra and its
Applications 353 (2002), 5--10 (with the history: D.~Rosenblatt 1957, J.~C.~Holladay and R.~S.~Varga 1958); A.~L.~Dulmage and
N.~S.~Mendelsohn, \emph{Gaps in the exponent set of primitive matrices}, Illinois Journal of Mathematics 8 (1964), 642--656 (which
exponents below $(n-1)^2+1$ never occur).
\item A.~A.~Markov, \emph{Extension of the law of large numbers to dependent quantities} (in Russian), Izvestiya
Fiziko-Matematicheskogo Obshchestva Kazan 15 (1906), 135--156.
\item P.~H.~Leslie, \emph{On the use of matrices in certain population mathematics}, Biometrika 33 (1945), 183--212.
\item S.~Brin and L.~Page, \emph{The anatomy of a large-scale hypertextual web search engine}, Computer Networks and ISDN Systems
30 (1998), 107--117; K.~Bryan and T.~Leise, \emph{The \$25{,}000{,}000{,}000 eigenvector: the linear algebra behind Google},
SIAM Review 48 (2006), 569--581.
\item E.~Seneta, \emph{Non-negative Matrices and Markov Chains}, 2nd ed., Springer, 1981; C.~D.~Meyer, \emph{Matrix Analysis and
Applied Linear Algebra}, SIAM, 2000, Chapter~8 (with a history of Perron--Frobenius theory).
\end{itemize}
